\begin{afterword}
\summaryhead

The post adds a fourth item to a series of virtues set out in the posts before it. After ``I want to become stronger,'' ``make a desperate effort'' and ``make an extraordinary effort'' comes ``shut up and do the impossible.'' ``Impossible'' here means that no one sees a way, not that there is a proof. What sets the fourth apart is its goal: not to improve, to try hard, or to try in new ways, but to win.

The illustration is the AI-Box Experiment. People who said no argument could make them release an AI played the gatekeeper, the author played the AI, and in the first two games both gatekeepers let the AI out. By the rules, the conversations stayed secret. The author says there was no trick: ``I just did it the hard way.''

The attitude the post asks for is to intend to solve the problem while keeping in view every reason it cannot be solved, expecting neither to win nor to lose, and refusing both cheap solutions and the comfort of faith. The author cannot put a probability on such efforts. In this sense, the author says, Friendly AI is impossible, so a proposal that needs only ordinary effort deserves suspicion. The post ends by reporting three later games, one won and two lost, and warns that the method is costly, that you can lose, and that it should be reserved for very special occasions.

\responsehead

The post is advice about an attitude, with one illustration. It includes its own warnings: the method is costly, you can lose, and it is for very special occasions. It also reports the author's two losses in later games.

The illustration is partly on record. The post's account of how the first two games began matches the 2002 mailing-list archive, and both gatekeepers confirmed in signed messages that they let the AI out. What happened in the games is secret by the rules, and the page itself calls the results ``strictly anecdotal evidence,'' as the notes on ``Einstein's Superpowers'' report. The three later games are recorded only in this post. So the example shows, on this record, that people sure no argument could move them were moved. It does not show how, so it illustrates the attitude only through the author's own account: ``I just did it the hard way.'' The post also concedes that the AI-Box game never seemed impossible to the author.

The advice depends on claims about minds that the post states without evidence: that believing you will fail lowers your standard, and that you can hold full intent to solve a problem while seeing every reason it is impossible. The first fits research on expectations and effort. In that research, though, people who weighed a desired future against present obstacles committed strongly only when they expected success.

The advice is applied to AI once: researchers who call Friendly AI impossible ``haven't even tried for the sake of trying,'' the author is ``pretty sure''; no example is given.

In short: advice to set out to win at problems that look impossible, with its own warnings, based on the author's experience and illustrated by games whose outcomes are on record but whose content is secret.
\end{afterword}
